\section{Introducción}
\label{sec:introduccion}
% Origen: 01-definicion-del-caso.md §2 y §4, 03-informe-tpi.md §1 y 10-antecedentes.md.

Constitución es la estación más cargada del subte de Buenos Aires. En junio de 2026 sus molinetes
registraron \num{1293543} ingresos, 2,3 veces los de la segunda estación de la red (Congreso de Tucumán,
Línea~D). Además, las 14 ventanas de 15 minutos más cargadas de las seis líneas fueron de Constitución
\parencite{sbase2026molinetes}. La razón es su ubicación: la estación de la Línea~C está debajo de la
terminal Plaza Constitución del Ferrocarril Roca. En esa terminal se vende el 29,3~\% de los pasajes del
Roca: 35,3 millones en 2024 \parencite[pp.~15 y 31]{cnrt2025amba}.

Este trabajo estudia el ingreso al subte por la batería de molinetes del vestíbulo Principal de esa
estación. El vestíbulo Principal da al hall ferroviario y concentra el 84,7~\% de los ingresos de la
estación. El estudio cubre los días hábiles entre las 07:00 y las 09:30. La pregunta es operativa y de
bajo costo: cuánto cambia la espera de los pasajeros si en el pico de la mañana se habilitan todos los
molinetes del vestíbulo en lugar de los que se observan en uso.

\subsection{Por qué el caso requiere simulación}

Con los promedios, el problema parece no existir. En la ventana más cargada (08:15, promedio de marzo a
junio de 2026) ingresan \num{2214} pasajeros en 15 minutos. Hay unos 19,5 molinetes con tráfico por
ventana. Eso da 7,6 pasajeros por minuto y por molinete, uno cada 7,9 segundos. Una validación dura del
orden de 2 a 3 segundos \parencite{tcqsm2013,seriani2025,wang2024}. Con una media de 2,57~s, la
utilización por molinete es $\rho \approx 0{,}32$. Un modelo M/M/$c$ con $c = 20$ servidores y esa carga
da una probabilidad de esperar del orden de $10^{-5}$: según la fórmula, no hay cola.

Sin embargo, en Constitución la cola existe y se ve todos los días. Los pasajeros no llegan repartidos en
el tiempo, sino en oleadas. En la franja estudiada arriban 47 trenes del Roca, uno cada tres minutos en
mediana, y cada uno descarga cientos de personas que caminan juntas hacia los molinetes
\parencite{trenes2026roca}. Después de cada arribo se forma una cola transitoria. Entre dos trenes, la
capacidad queda en gran parte ociosa. Es un caso de \emph{congestión transitoria por arribos en lote sobre
un sistema subutilizado en promedio}. Ahí los modelos analíticos estacionarios fallan, porque promedian
justamente lo que se quiere medir. Por eso se usa simulación de eventos discretos.

El tiempo que la oleada tarda en llegar a los molinetes también importa. Ese tiempo depende de la
caminata por el andén y el hall, y de cómo se reparten las personas entre las colas. Por eso el modelo
principal se construyó con la Pedestrian Library de AnyLogic \parencite{anylogic2026help}: cada pasajero
camina por un vestíbulo, elige un molinete y espera en la cola de ese molinete.

\subsection{Antecedentes}

Varios estudios recientes simularon estaciones de metro con AnyLogic y midieron parámetros que este
trabajo usa como referencia.
\begin{itemize}
  \item \textcite{seriani2025} observaron los molinetes de una estación del Metro de Valparaíso y midieron
        un tiempo de paso medio de 2~s. Registraron además un sesgo fuerte por cercanía: el molinete
        más próximo al recorrido es el más usado. Validaron su modelo comparando el uso simulado de cada
        molinete con el observado.
  \item \textcite{wang2024} midieron con video 813 pasos por molinete en Pekín, con medias de 2,91 a 3,49~s
        según el medio de pago.
  \item \textcite{tian2026} modelaron una estación de trasbordo que recibe pasajeros de un ferrocarril, la
        situación análoga a la del Roca y la Línea~C. Encontraron que el 88,2~\% usa el grupo de molinetes
        más cercano a la salida del tren.
  \item \textcite{li2024} calibraron con video el modelo de fuerza social de AnyLogic y obtuvieron una
        velocidad deseada de 1,37~m/s, dentro del rango de 1,1 a 1,5~m/s de la literatura.
  \item \textcite{zhen2024} usaron datos abiertos de molinetes de Nueva York, como aquí los de SBASE, pero
        con arribos de Poisson y sin oleadas.
\end{itemize}
Lo que agrega este trabajo es representar las oleadas a partir del horario real de los trenes y
distribuir entre ellos la demanda observada por ventana.

\subsection{Organización del informe}

El informe sigue los diez pasos de un estudio de simulación de la consigna de la cátedra
\parencite{utn2026consigna,law2015}. La sección~\ref{sec:paso1-formulacion} formula el problema y los
escenarios. La sección~\ref{sec:paso2-datos-modelo-conceptual} describe los datos y el modelo conceptual.
Las secciones~\ref{sec:paso3-validacion-conceptual} a~\ref{sec:paso6-validacion} cubren la validación
conceptual, la construcción y verificación del programa, los pilotos y la validación del modelo
programado. Las secciones~\ref{sec:paso7-diseno-experimentos} a~\ref{sec:paso10-conclusiones} presentan el
diseño de experimentos, las corridas, el análisis estadístico y las conclusiones.

Un aviso vale para todo el documento. El grupo reside en Rosario y no pudo medir en Constitución. Los
parámetros que no tienen dato publicado se tomaron de fuentes técnicas citables, como rangos con un valor
intermedio, y así se declaran. El plano del vestíbulo es hipotético hasta recibir el oficial, que se
solicitó a SBASE. Las conclusiones sobre la espera valen dentro de esos supuestos.
